\section{AI-assisted proofs and a proposed reporting standard}\label{sec:standard}

This section describes current practice in the use of AI for mathematical proofs, compares the approach of this
paper with recent large-scale efforts, and proposes the reporting standard. The standard is applied to the case
study in Section~\ref{sec:collab}.

\subsection{Current practice}\label{sec:std-practice}

In most papers in our field that acknowledge AI, a general-purpose chatbot was asked for help with a step, a lemma,
a computation or the wording, and the acknowledgement says so in a sentence. The proofs remain the author's, and a
one-line declaration is adequate. A different mode of work has appeared in which AI agents, instances of a language
model given a role, tools and a task, are organized to propose, prove, criticize and formalize results, and the
human contribution is to choose problems, direct the work and decide what to accept.

The most visible examples are the releases of OpenAI in 2026, which illustrate several distinct ways of combining
AI discovery with verification. In February, OpenAI submitted proof attempts by an internal model to First Proof,
a set of ten research-level problems posed by experts in their fields \citep{openai_firstproof_2026}. The model
was run ``with limited human supervision'': the researchers sometimes suggested retrying strategies that had
appeared fruitful, asked the model to expand or clarify parts of a proof after expert feedback, facilitated a
back-and-forth between the model and ChatGPT ``for verification, formatting, and style'', and for some problems
presented the best of a few attempts selected by human judgment. One attempt first believed likely correct was
later judged incorrect in the light of the official commentary, OpenAI states that the process was not as clean as a
properly controlled evaluation, and the accompanying preprint includes an appendix of prompt patterns. In
August, OpenAI published ten results in mathematics and theoretical computer science, with certificates in the
proof assistant Lean \citep{openai_tenproofs_2026}. In September it announced a resolution of the Navier--Stokes
problem, with a written proof and a Lean formalization \citep{openai_navierstokes_2026}. OpenAI describes a
system of coordinating agents, powered by an internal model, with access to a cached copy of the internet and to
code execution, and subdivided into groups that communicate internally; the group that produced the result
involved of the order of ten thousand concurrent agents. Different groups were prompted with different variants of
the problem statement, two leading to a proof of regularity and two to a counterexample. A run on easier problems
produced, with nearly a hundred agents over about fifty hours, a result on the Euler equations, after which the
researchers moved agents to Navier--Stokes, prompted them with the Euler result, and switched them to a further
trained version of the model. Groups were encouraged to explore diverse approaches and were then
``cross-pollinated'' by using a coding agent to consolidate the most useful insights of each group into follow-up
prompts. The resolution came after about 88 hours and 2.7 million messages between agents, and the Lean
formalization took another 17 hours. In October OpenAI released about seven hundred
manuscripts produced by an internal model by one fixed procedure on some four thousand posed problems, together
with a Lean library covering a part of the results and a change log recording retractions and repairs
\citep{openai_math_2026}; this release did not use the large communicating configuration of the Navier--Stokes
run. The repository describes the manuscripts as ``raw research data, rather than polished academic papers'', and
states that unformalized results could contain errors. Later in October, OpenAI announced three further results, in
harmonic analysis, complexity theory and quantum field theory, again described its manuscripts as raw research
data, and tested three ways of turning them into understanding: releasing the materials and leaving researchers to
organize themselves; forming an external group of collaborators while releasing the materials; and forming an
external group and leaving the release to it \citep{openai_academic_2026}. The collaborations follow four stated
principles: transparency, no academic credit claimed by OpenAI for discoveries of its models, recognition of the
groups for the understanding they develop, and independence of the groups. OpenAI notes that the choice of
coordinators gives some researchers earlier access to the materials than others.

The mathematical community has responded on three points that bear directly on reporting. First, on
\emph{understanding and disclosure}: the Advisory Group on Mathematics and Artificial Intelligence (AGMAI), a group
of mathematicians independent of AI companies, distinguishes papers that a human understands, to which the usual
norms apply, from results not yet understood by anybody, and asks for the latter that the model, the prompts, a
summary of the reasoning, the time and the cost of computation be disclosed, that it be stated how problems were
chosen and how many were not solved, and that results be deposited in repositories not controlled by the releasing
company \citep{agmai_guidelines_2026}. Second, on \emph{formal verification}: a formal check establishes the formal
statement, not that it says what the paper says. A human audit of the August release found no confirmed substantive
error in a principal result but stressed that Lean ``checks the encoded proposition'' while human scrutiny must
establish its intended meaning, and recommended a standardized report header \citep{sienicki_audit_2026}; a study
of the Navier--Stokes release argues that its formal proof does not correspond to the natural-language proof
\citep{bastounis_lost_2026}. The Lean community's own checklist for formal claims asks the same questions: whether
the proof is complete, which axioms it depends on, and whether it proves the statement that is claimed
\citep{lean_didyouproveit}. Third, on \emph{records}: press reports noted that only a small fraction of the
October manuscripts came with summaries of the model's reasoning, and that natural-language and formal artifacts
were not linked by machine-readable metadata \citep{fernholz_techcrunch_2026,howlett_sciam_2026}.

\begin{discussion}[the prevailing practice and the role of repositories]
In preparing this paper the author observed that most reported uses of AI in our field are of the first kind, a
chatbot that helps to establish a lemma and is then acknowledged, and that the second kind needs a fuller account:
the form of the panel, the use of a proof assistant, and the place of the human in the loop. The author and the
main assistant also noted that this way of working resembles programming more than writing: the proof assistant
code, the records of the process and the manuscript source are kept and published together in a version-controlled
repository, rather than as a single preprint. These observations motivated the items on workflow, formal
verification and records in Table~\ref{tab:standard}.
\end{discussion}

\subsection{The approach of this paper compared with large-scale efforts}\label{sec:std-compare}

The case study used a panel of AI agents run by one researcher (Section~\ref{sec:collab-org}). Its design shares
several features with the large-scale efforts, and differs from them in scale and in the role of the human.
Table~\ref{tab:compare} summarizes the comparison. OpenAI's descriptions cover the organization of the agents into groups, the
variants they received and the consolidation of their findings, but not roles such as reviewer or judge within a
group, and the table does not attribute such roles to them.

\begin{table}[ht]
\centering\footnotesize
\begin{tabularx}{\textwidth}{@{}>{\raggedright\arraybackslash}p{2.6cm}>{\raggedright\arraybackslash}X>{\raggedright\arraybackslash}X@{}}
\toprule
\textbf{Feature} & \textbf{Large-scale efforts (OpenAI, 2026)} & \textbf{This paper}\\
\midrule
Parallel work on variants & communicating groups prompted with different variants of the problem statement
  (Navier--Stokes) & proposal agents on distinct angles of one question\\
Synthesis & a coding agent consolidates the most useful insights of each group into follow-up prompts & judge agents
  rank the proposals, choose one and graft parts of others onto it\\
Criticism & model-to-model checking and expert feedback, with human selection among attempts (First Proof);
  dedicated adversarial reviewers not documented & each proposal and proof attacked by agents instructed to refute
  it; a review by a model of a second provider\\
Formal verification & Lean certificates for part of the results & Lean check of the loss-ratio-specific steps of
  the first formulation, in an abstract setting (Section~\ref{sec:collab-lean})\\
Release & public repositories with change logs of retractions and repairs & public repository with version
  history, proof records and process records\\
Error record & retractions and repairs in a change log & every error found in AI-generated mathematics, with how it
  was found (Table~\ref{tab:errors})\\
Scale & of the order of $10^4$ concurrent agents, 88 hours and 2.7 million messages for one problem; thousands
  of posed problems in the October release & a few to about ten agents per stage; one problem\\
Configurations & a large communicating run for one problem, with the model updated during the run; a fixed
  procedure for thousands of problems & one panel design, reused with new roles for the revision; several model
  versions across sessions (declaration)\\
Role of the human & choice of problems and of where to direct computation; suggestions of strategies, requests
  for clarification and selection among attempts (First Proof); external working groups after release
  & choice of the problem and of every formulation; objections to AI drafts changed the result\\
Credit and responsibility & no academic credit claimed for discoveries of the models; credit to external groups
  for the understanding they develop & the author takes responsibility for the result; the division of labour
  records what the AI contributed\\
Disclosure of prompts & prompt patterns in an appendix (First Proof); for Navier--Stokes, offered privately to
  concurrent researchers & the prompts of the latest panel are kept with its records\\
\bottomrule
\end{tabularx}
\caption{The approach of the case study compared with large-scale multi-agent efforts in mathematics.}\label{tab:compare}
\end{table}

The resemblance is in the design principles: several independent attempts from different angles, selection and
synthesis of the most useful findings (by a coding agent in the Navier--Stokes run, by judge agents here),
criticism, formal verification of what can be formalized, and release of the records with
a history of corrections. These principles do not depend on scale, and a single researcher can apply them to a
problem in risk management at modest cost. The closest analogue to the case study is OpenAI's First Proof work,
in which humans suggested strategies, asked for clarification after expert feedback, used a second model to check
the proofs and selected among attempts, which is the role the author and the review with ChatGPT played here; what
the case study adds is a record of the errors found at each step and of the author's interventions. The difference is in the human role. In the large-scale efforts the
human chooses problems and allocates computation, and the results may not yet be understood by anybody; in the case
study the author directed every stage, and the central idea of the result came from the author's objections to
drafts produced by the agents. The assignment of credit differs accordingly: OpenAI claims no academic credit for its models' discoveries and
credits external groups with the understanding they develop, whereas here the author, who directed the work and
understands the result, takes responsibility for it, and the report records what came from the AI. The paper
therefore falls under the first category of the AGMAI guidelines, a result
that its author understands, but its proofs were developed by agents, which is why the fuller report of
Section~\ref{sec:collab} is given.

\begin{discussion}[the design of the panel]
When the author decided to make the use of AI central to the paper, the panel was described to the main assistant
as organized in the manner of the recent multi-agent efforts in mathematics, and the paper was conceived as having
two parts: the mathematical question, and the demonstration, open to inspection, of how AI was used to answer it.
To check this description, two accounts of OpenAI's methods were prepared independently: one by an agent of the
main assistant, which read OpenAI's public repositories, the arXiv audits, the AGMAI guidelines and press reports
but could not retrieve OpenAI's own pages, and one by ChatGPT, at the author's request, which cited OpenAI's posts
directly. The author then retrieved OpenAI's Navier--Stokes and First Proof posts and its October post on academic
collaborations in a browser, and the main assistant
compared them with both accounts; the posts confirm the facts that both accounts reported, and the text above
quotes them directly. The two accounts agree on
the facts used here. Both find that the resemblance holds for the principles
listed above, that OpenAI used model-to-model checking and coordinated groups of agents, and that no published
record specifies dedicated reviewer or judge roles; the comparison in Table~\ref{tab:compare} is stated
accordingly. The use of a second provider to cross-check a description of the literature is an instance of the
guardrail of Section~\ref{sec:collab-guard}.
\end{discussion}

\subsection{The proposed standard}\label{sec:collab-standard}

We propose that a mathematical paper in risk management or actuarial science whose results were obtained with
substantial AI assistance report the items in Table~\ref{tab:standard}, in the tier that fits its use of AI
(Section~\ref{sec:std-tiers}). For each item the table states whether an editor or reviewer can verify it from
public artifacts (V), can inspect it without being able to establish its completeness (P), or must take it on the
author's word (S), and where it is reported for the case study.

\begin{table}[ht]
\centering\footnotesize
\begin{tabularx}{\textwidth}{@{}>{\raggedright\arraybackslash}p{2.5cm}>{\raggedright\arraybackslash}X>{\centering\arraybackslash}p{0.75cm}>{\centering\arraybackslash}p{0.55cm}>{\raggedright\arraybackslash}p{2.5cm}@{}}
\toprule
\textbf{Item} & \textbf{What to report} & \textbf{Tier} & \textbf{Ver.} & \textbf{Case study}\\
\midrule
1. Tools & providers, models and versions, mode of access, dates of use & 1--3 & S/P & declaration\\
2. Division of labour & who posed the question, chose the formulation and decided on content; which tasks were
  given to AI & 1--3 & S & Secs.~\ref{sec:intro}, \ref{sec:collab-org}\\
3. Workflow & the agents or sessions used, their roles and order, the review loops, and the points at which the
  author intervened & 3 & P & Sec.~\ref{sec:collab-org}, \ref{sec:collab-revision}\\
4. Knowledge base & how the literature was gathered, how reliable each statement was taken to be, which cited
  statements were checked against the original works & 3 & P & Sec.~\ref{sec:collab-kb}\\
5. Error log & every error in AI-generated mathematics that was found, how it was found, how it was resolved
  & 2--3 & S/P & Table~\ref{tab:errors}\\
6. Guardrails & the checks applied: adversarial review, review by a second system, recomputation of numbers,
  counterexamples, self-contained proofs of uncertain citations & 2--3 & P & Sec.~\ref{sec:collab-guard}\\
7. Formal verification & which results were checked by a proof assistant, in what setting, which steps were taken
  as hypotheses, and how the formal statements were compared with the paper & 3\textsuperscript{*} & V &
  Sec.~\ref{sec:collab-lean}\\
8. Verification status & for each part of the paper, whether it was reviewed by AI, by a human, or formally
  checked, and what has not been checked & 2--3 & S/V & Sec.~\ref{sec:collab-status}\\
9. Records & the formal proofs, the scripts and the records of the AI process, in a public repository with version
  history & 3\textsuperscript{*} & V & \texttt{ai\_process/}, \texttt{skill/}, \texttt{agents/} in the
  repository\\
\bottomrule
\end{tabularx}
\caption{Proposed items for reporting AI-assisted proofs. Tier: the tiers in which the item applies
(Section~\ref{sec:std-tiers}); $^*$ encouraged, not required. Ver.: V verifiable from public artifacts, P open to
inspection, S self-declared (Section~\ref{sec:std-verify}).}\label{tab:standard}
\end{table}

\subsection{Proportionality}\label{sec:std-tiers}

A reporting requirement must be proportionate to the use of AI, or it will be ignored or will burden papers that do
not need it. We propose three tiers.
\begin{enumerate}[label=(\roman*)]
\item \emph{Incidental use}: AI was used for language, literature search, coding unrelated to proofs, or similar
  tasks, and contributed no mathematical content. Items~1 and~2, which a declaration of AI use covers.
\item \emph{Assisted proofs}: AI proposed or checked individual steps, lemmas, examples or counterexamples, which
  the author verified and which appear in the paper. In addition, items~5, 6 and~8: the errors found in the AI
  contributions, the checks applied to them, and the verification status of the parts concerned.
\item \emph{Agent-developed proofs}: proofs of main results were drafted or developed by AI agents, with the
  author directing and reviewing. All nine items, with formal verification and a public repository encouraged.
\end{enumerate}
The tier is determined by the most substantial use: a paper whose main theorem was first proved by agents is in
tier~(iii) even if most of its text was written by the author. The case study is in tier~(iii).

\subsection{Verifiability and incentives}\label{sec:std-verify}

A journal cannot enforce rules that it cannot check; it can, for instance, not prohibit the use of AI in proofs.
The standard therefore asks for information rather than imposing restrictions, and it relies on two mechanisms to
make honest reporting the author's best course.

First, items differ in what can be checked. A formal proof can be rebuilt by anyone (V), and a repository with
version history shows when each part of the paper and of the formal proof was written and changed (V). Process
records and scripts can be inspected and rerun, although their completeness cannot be proved (P). The division of
labour and the error log are statements of the author (S). Self-declared items are not unverifiable in isolation:
a version history, records and a formal proof constrain them, and an inconsistency between the report and the
records is visible to any reader.

Second, the items that cost the author effort are those that save the reviewer effort. A formal check removes the
verified steps from the reviewer's work, a verification-status table directs the reviewer to what remains, and an
error log shows where AI drafts went wrong and how this was detected. A reviewer who is given these can review an
AI-heavy paper faster and with more confidence than one who is told only that AI was used. An author who reports
errors found and corrected signals that the result was checked; an author who reports none in an agent-developed
paper invites the question of how hard the result was examined. Under these conditions a candid report is in the
author's interest.

\subsection{Formal verification, repositories and the reviewer}\label{sec:std-repo}

Formal verification in a proof assistant should be encouraged for AI-assisted proofs, but not required. Mathlib
does not yet cover all the mathematics used in risk management, and in the case study the classical representation
step could not be formalized and entered as a hypothesis (Section~\ref{sec:collab-lean}). When a formal check is
made, its scope must be reported precisely, as the Navier--Stokes discussion shows: what was formalized, in what
setting, which steps were hypotheses, and how the formal statements were compared with the statements of the paper.

A public repository that holds the formal proofs, the scripts and the records of the process, with a version
history, should eventually be expected for agent-developed work. Such a repository changes the reviewer's task:
part of the verification is done by the proof assistant, and the reviewer's attention shifts to whether the formal
statements and definitions say what the paper says, whether the hypotheses are reasonable, and whether the
unverified parts are sound. This requires skills that the actuarial research community will have to acquire,
notably reading formal statements and rebuilding a formal proof, and it raises the question of how reviewers may
themselves use AI. The standard does not settle these questions.

\begin{discussion}[the editorial perspective]
An editor of an actuarial journal, asked for an opinion on an earlier version of the standard, pointed out that
editors are at present more concerned with the volume of submissions and the strain on reviewers than with
correctness, that a reporting framework should apply only to AI-heavy work and be proportionate and flexible, that
rules which cannot be verified should not be imposed and that a framework should give authors an incentive to
report honestly, that formal verification should be encouraged although requiring it is premature, and that
repositories should eventually be expected and will change the role of reviewers. In reading these comments with
the main assistant, the author identified the changes made in this version: the tiers of
Section~\ref{sec:std-tiers}, the verifiability column of Table~\ref{tab:standard} and
Section~\ref{sec:std-verify}, the emphasis on the reviewer's work in this subsection, and the restructuring of the
paper around the standard, with the mathematical result as its case study and its proofs in the Supplement.
\end{discussion}

\subsection{Using the standard: recommendations and a reusable skill}\label{sec:std-use}

The experience of the case study suggests the following recommendations.

\emph{For authors who apply the standard.} Decide the tier at the start of the work rather than at submission, and
set up the repository then. Record errors, discussions and decisions when they occur; reconstructing them
afterwards is unreliable. Keep the prompts and outputs of every panel. If a formal check is planned, formalize the
definitions early, since this exposes ambiguities in the statements. State the verification status part by part and
with dates, so that a later revision of the mathematics does not silently fall outside an earlier check. Name readers
or referees whose comments changed the work only with their consent.

\emph{For researchers who want to carry out a similar case study.} Choose a problem whose answer is not known but
whose correctness can be checked, and list its degenerate cases in advance; most of the errors found here occurred
at such cases. Give the proposal agents distinct angles, one of which should always be the limits of what can be
achieved. Give each proposal its own adversarial checker, and test claims at the points where they can fail rather
than at random points (Table~\ref{tab:errors}). Use a model of a second provider for review. Above all, keep the
human role active: in the case study the decisive steps came from the author's objections, not from the agents.

\emph{For editors and reviewers.} Read the verification status first, since it shows where the unverified parts are.
Where a formal proof is provided, rebuild it and spend the review effort on whether the formal definitions and
statements say what the paper says. Compare the report with the version history of the repository, and check a few
entries of the error log against the records.

\emph{Practical advice for proofs with agents.} The following points summarize what the case study taught about
working with agents on proofs; the last three are general experience rather than findings of this project.
\begin{itemize}[nosep]
\item \emph{Fix the statement before proving it}, with all quantifiers, the domain and the conventions, and ask each
  agent to restate it; differences between the restatements expose ambiguities.
\item \emph{Give every agent the list of degenerate cases} and of established facts; almost every error in
  Table~\ref{tab:errors} occurred at a degenerate case, and a list of known facts prevents their rediscovery.
\item \emph{Freeze the statement while a proof is checked.} The typical failure is a correct argument for a slightly
  weaker statement; checkers should compare what was proved with what was claimed.
\item \emph{Search for counterexamples before proofs.} The witnesses of Section~\ref{sec:necessity} and
  Proposition~\ref{prop:nofinite} arose in this way.
\item \emph{Ask for the dependency structure of a proof}, with the hypotheses of each lemma, so that circular
  citations become visible; treat ``clearly'' as a gap to be expanded.
\item \emph{Require concrete objections}, a counterexample or a failing step, and test on small finite models in
  exact arithmetic at the points where a claim can fail.
\item \emph{Check citations that carry weight} against the source or reprove them; agents misremember hypotheses
  and numbering.
\item \emph{Formalize definitions and statements first}, and have a human compare them with the paper before formal
  proofs are attempted.
\item \emph{Keep the failed approaches}; several became impossibility results of the case study.
\item \emph{Ask what a formulation assumes.} The first complete theorem of the case study supplied the premium and
  the risk measure by assumption, and the author's objection to it produced the central idea of the paper.
\item \emph{Give checkers the statement and the proof, not the prover's reasoning}, so that they are not anchored on
  it; stop checking only when consecutive rounds find nothing new; and keep the working context in files rather
  than in one long session.
\end{itemize}

\emph{A reusable skill.} To make the method usable by others, the protocol of the case study has been written as a
\emph{skill}: a structured instruction file that an AI assistant loads when a task matches its description, and that
a human can read as a written protocol. It is in the directory \texttt{skill/} of the repository. It sets out the
roles of author, assistant and agents; the repository layout; the stages of the work (knowledge base, proposal
panel, adversarial check, judge, checking runs and repair, second provider, numerical checks, formal verification);
the pitfalls met in the case study; templates for the error log, the records of discussions and revisions, and the
verification status; the nine items with their tiers; and notes on running the panels in different tools, with the
workflow script of the revision as an example. Assistants that support skill files load it directly; with other
tools, the roles can be run as separate sessions that share one context document. The skill encodes procedure, not
judgment, and it was distilled from one case study; it is offered, like the standard, as a starting point for others
to adapt.

\emph{Reusable agents.} The agent roles of the case study are also provided in a form that others can use, in the
directory \texttt{agents/} of the repository. Each agent is a text file with a short header, giving its name, a
description of when it should be used and the tools it may use, followed by its instructions; this is the format in
which Claude Code defines subagents, and with other assistants the instructions can be given to a separate session
as its role. The roles are those of Figures~\ref{fig:workflow} and~\ref{fig:panel} and of the revision panel: a
literature agent for one strand, a proposer working from an assigned angle, an adversarial checker instructed to
refute one proposal, a judge, a chair who writes the note to the author, a prover, a checking run with a specific
brief, a Lean formalizer, and a prompt for an independent review by a model of another provider. All of them read
a shared context file, for which a template is given, that fixes the setting, the target, the established facts,
the degenerate cases, the author's constraints and the notation; in the case study, a precise context of this kind
was the single most effective safeguard. A workflow script runs a complete panel, with proposers in parallel, each
followed by its checker, then the judge and the chair. The instructions of each role encode the lessons listed
above; for instance, the checker compares the proved statement with the claimed one and tests claims at the points
where they can fail. The agent files and the workflow were tested on a problem with a known answer, the
characterization of the functionals $a\max+(1-a)\min$ on a two-point space by cash additivity, comonotonic
additivity, monotonicity and law invariance. The panel of six agents recovered the answer with an independence
witness for each axiom; an adversarial checker found three presentational errors in one proposal, and the chair
marked as a gap the one step that needs the axiom of choice. The test also showed that agent definitions are read
when a session starts, so the workflow can instead tell each agent to read its role file directly. The test run is
included in \texttt{agents/} as a worked example.

\begin{discussion}[a skill for others]
Towards the end of the work the author asked that the method be written up as a skill file, so that the paper could
be used by others, and asked whether such a file can serve someone who wants to carry out a similar case study. The
author also asked for the agents designed for the project to be written up in a form that others can use, which led
to the agent files described above. The conclusion was that both can serve others, as a protocol, a set of roles and
a set of templates: the file is plain text, it is loaded directly by
assistants that support skills and can be followed by hand otherwise, and it records the pitfalls that the case
study met. It cannot transfer the author's understanding of a problem, which the case study shows to be decisive,
and it should be revised as it is used on other problems.
\end{discussion}

\subsection{Limits of the proposal}

The list is a first proposal. It is based on one project, and it will need to be tested on other papers, adapted
to results whose proofs are checked formally in full, and revised as the tools change; what AI-assisted research
looks like may change within months. It addresses the reporting of results whose authors understand them, and it
does not address results produced at scale without human understanding, for which the AGMAI guidelines are more
appropriate. Agreement among AI reviews does not establish correctness, and a formal check covers only what was
formalized, so the verification status must be stated part by part.
